\documentclass{article}
\usepackage{graphicx} % Required for inserting images
\usepackage{amsmath}
  \RequirePackage[french]{babel}
  \RequirePackage[T1]{fontenc}
  \RequirePackage[utf8]{inputenc}

\title{STT-2902 - Modélisation statistique}
\author{Julien Miron}


\begin{document}

\section*{Équipe 3 – Intervalles de confiance}

\textbf{Session :} Automne 2026

\textbf{Style imposé : atelier expérimental.} Votre présentation doit prendre la forme d'un atelier pratique où les autres étudiants manipulent des données réelles ou simulées et construisent eux-mêmes des intervalles de confiance. Ce style est obligatoire : l'essentiel du temps doit être consacré à l'expérimentation par l'auditoire, pas à un exposé.

\vspace{0.5cm}
\textbf{Objectif :} Faire comprendre, par l'expérimentation, ce qu'un intervalle de confiance représente réellement et comment son niveau de confiance se manifeste sur plusieurs échantillons.

\vspace{0.5cm}
\textbf{Déroulement imposé} :
\begin{itemize}
    \item Une courte mise en contexte (au plus 10 minutes) présentant l'expérience et les formules nécessaires;
    \item Une activité d'échantillonnage concrète (ex. : jetons numérotés, dés, données distribuées) où chaque équipe tire son propre échantillon et calcule son intervalle de confiance;
    \item Une mise en commun des résultats de toutes les équipes, avec une visualisation de l'ensemble des intervalles obtenus;
    \item Un retour guidé sur ce que l'expérience illustre.
\end{itemize}

\vspace{0.5cm}
\textbf{Requis} :
\begin{enumerate}
    \item Chaque équipe de l'auditoire doit construire au moins un intervalle de confiance à partir de son propre échantillon;
    \item Vous devez clairement expliquer ce qu'un intervalle de confiance de niveau $(1-\alpha)$ nous donne comme information;
    \item Produire un graphique montrant les intervalles de confiance de toutes les équipes et la vraie valeur du paramètre. Vous pouvez le produire au tableau, aucun logiciel n'est obligatoire;
    \item Discuter de l'effet de la taille d'échantillon et du niveau de confiance sur la largeur de l'intervalle.
\end{enumerate}

\end{document}
